\clearpage
\chapter*{\texthebrew{תקציר}}
\addcontentsline{toc}{chapter}{Abstract (Hebrew)}
\markboth{Abstract (Hebrew)}{Abstract (Hebrew)}

\begin{otherlanguage}{hebrew}
\begin{flushright}
\label{sec:abstract-hebrew}

\textbf{ארכיטקטורות למידה עמוקה לתקשורת ממסר דו-קפיצתית: מחקר השוואתי של אסטרטגיות ממסר קלאסיות ומבוססות רשתות נוירונים}

חיבור זה מציג מחקר השוואתי של אסטרטגיות ממסר (relay) קלאסיות ומבוססות רשתות נוירונים עבור מערכות תקשורת שיתופית דו-קפיצתית. תשע שיטות ממסר נבדקות: שתי גישות קלאסיות — הגברה-והעברה (AF) ופענוח-והעברה (DF) — ושבע שיטות מבוססות למידה עמוקה: למידה מפוקחת (MLP מינימלי, ממסר Hybrid המתאים את עצמו לרמת ה-SNR), מודלים גנרטיביים (VAE, CGAN עם אימון WGAN-GP), וארכיטקטורות רצפים (Transformer עם קשב-עצמי רב-ראשי, Mamba S6 עם מרחב-מצבים סלקטיבי, Mamba-2 SSD עם כפל מטריצות semi-separable מקבילי).

ההערכה מבוצעת על פני שישה תצורות ערוץ: AWGN, דעיכת Rayleigh, ודעיכת Rician ($K=3$) בטופולוגיית SISO, וכן MIMO $2\times2$ עם שלוש שיטות איזון — ZF, MMSE, ו-SIC. הניסויים משתמשים באפנון BPSK עם סימולציית מונטה קרלו (100,000 ביטים לכל נקודת SNR) ורווחי סמך של 95%. ניסויי הרחבה מכסים QPSK, 16-QAM ו-16-PSK.

הממצאים המרכזיים: ראשית, ממסרים מבוססי רשתות נוירונים מראים יתרון סלקטיבי בתחום SNR נמוך ($0$–$4$ dB), בעיקר בערוצי AWGN ו-Rician ובחלק מתצורות ה-MIMO, אך יתרון זה אינו אוניברסלי. שנית, ממסר DF הקלאסי שולט בתחום SNR בינוני-גבוה ($\ge6$ dB) עם אפס פרמטרים. שלישית, מחקר מורכבות מגלה יחס U-הפוך: רשת מינימלית בת 169 פרמטרים משתווה למודלים גדולים פי 100, בעוד שמודל בן 11,201 פרמטרים מציג התאמת-יתר. השוואה מנורמלת ל-3,000 פרמטרים מראה שהפער בין הארכיטקטורות מצטמצם בקנה מידה שווה, כאשר VAE הוא בעל הביצועים הנמוכים ביותר — ממצא המעיד שמספר הפרמטרים, ולא הארכיטקטורה, הוא הגורם המשפיע העיקרי. Mamba-2 SSD מתאמן מהר ב-35% מ-Mamba S6 תוך השגת BER זהה. במערכות MIMO, MMSE עולה על ZF ו-SIC מספק שיפור נוסף.

אסטרטגיית הפריסה המומלצת היא ממסר Hybrid המשלב עיבוד AI ב-SNR נמוך עם DF ב-SNR גבוה (169 פרמטרים, כ-0.7 KB, פחות מ-3 שניות אימון). עבור אפנונים מסדר גבוה, ניסוח הממסר כ\textbf{מסווג דו-ממדי משותף} על פני כל 16 נקודות הקונסטלציה מבטל את רצפת ה-BER המבנית: שלושת הוריאנטים המובילים (VAE, Transformer, MLP) משיגים BER קרוב לאפס ב-20 dB, ומשתווים לראשונה לביצועי DF. מחקר הזרקת CSI על פני 48 וריאנטים מגלה תלות-אפנון: CSI מזיק ל-16-QAM אך מועיל ל-16-PSK. הממצא המרכזי הוא שמורכבות המודל צריכה להיות מותאמת למורכבות המשימה — ארכיטקטורות מינימליות מספיקות לממסר BPSK/QPSK, ובחירת הפרדיגמה חשובה פחות מגודל מודל נכון ורגולריזציה מתאימה.

\end{flushright}
\end{otherlanguage}

